\documentclass[10pt,twocolumn]{article}

% ── Packages ──────────────────────────────────────────────────────────────────
\usepackage[a4paper, top=1.8cm, bottom=1.8cm, left=1.4cm, right=1.4cm, columnsep=0.5cm]{geometry}
\usepackage{times}
\usepackage{graphicx}
\usepackage{amsmath}
\usepackage{booktabs}
\usepackage{array}
\usepackage{hyperref}
\usepackage{cite}
\usepackage{enumitem}
\usepackage{caption}
\usepackage{microtype}
\usepackage{amssymb}
\usepackage{balance}

\hypersetup{
  colorlinks=true,
  linkcolor=black,
  citecolor=black,
  urlcolor=black
}

% ── Title ─────────────────────────────────────────────────────────────────────
\title{%
  \textbf{A Blockchain-Based Provenance System for Ensuring}\\
  \textbf{AI Training Data Integrity and Auditability}
}

\author{
  Xinyuan Hu\\
  \small Master of Software Engineering, Yoobee Colleges, Auckland, New Zealand
}

\date{}

% ── Document ──────────────────────────────────────────────────────────────────
\begin{document}

\maketitle
\thispagestyle{empty}

% ── Abstract ──────────────────────────────────────────────────────────────────
\begin{abstract}
The trustworthiness of industrial AI systems increasingly depends on the verifiable integrity
of their training process, yet existing solutions leave three critical engineering gaps
unaddressed: the absence of a deployable hybrid storage architecture, insufficient granularity
in dynamic process auditing, and the lack of atomic consistency guarantees between on-chain
records and off-chain data. This paper presents a blockchain-based AI training traceability
system that bridges all three gaps. The proposed system combines a four-layer architecture
integrating a public blockchain, IPFS, and OrbitDB with a non-intrusive Python SDK that
captures training metadata at Session, Epoch, and Step granularity through runtime
monkey-patching. A simplified two-phase atomic commit protocol, backed by a
compensation-log recovery mechanism, ensures idempotent fault recovery without requiring
changes to user training code. Controlled experiments demonstrate that the system restricts
overall training overhead to $+3.28\%$---strictly below the 5\% industrial tolerance
threshold---while achieving 100\% tamper detection across all adversarial scenarios, P99
one-click verification latency below 500\,ms, linear storage growth ($R^2 = 0.99985$),
and 100\% fault recovery with zero idempotency violations. The result is a directly
deployable engineering solution for compliance-oriented AI governance, with direct
applicability to frameworks such as the EU AI Act.
\end{abstract}

\textbf{Keywords:} blockchain; AI training traceability; IPFS; OrbitDB; atomic consistency;
hash chain; non-intrusive SDK; Design Science Research

% ── 1. Introduction ───────────────────────────────────────────────────────────
\section{Introduction}

With the large-scale deployment of industrial AI applications, the credibility of model
training, the compliance of data sources, and the auditability of the process have become
core requirements of algorithm governance~\cite{soldatos2021}. Regulatory frameworks such
as the EU AI Act mandate complete training data lineage, process audit logs, and model
version archives for high-risk AI systems operating in domains such as medical diagnosis,
financial credit scoring, and recruitment automation. Simultaneously, the reproducibility
crisis in AI research---where many published models cannot be independently reproduced due
to incomplete process records---has intensified demand for cryptographic-level evidence of
training provenance.

Despite this dual pressure from regulation and scientific rigour, existing solutions fall
short in three interconnected engineering dimensions. First, storing large-scale training
metadata directly on a public blockchain is neither cost-effective nor practically viable due
to gas costs and execution overhead~\cite{kurada2023}. Second, most auditing schemes focus
only on static dataset integrity and ignore the dynamic evolution of hyperparameters,
intermediate checkpoints, and per-step metrics throughout training~\cite{sciammarelli2026}.
Third, current solutions lack atomic commit, crash recovery, and log-compensation
mechanisms, leaving the traceability chain fragile under write interruptions or node
failures~\cite{belchior2022}. The intersection of these gaps means that even engineering
teams already using dedicated experiment tracking tools such as MLflow or DVC cannot
independently verify the integrity of their own records~\cite{idowu2024}.

This paper reports on a full-cycle design, implementation, and evaluation of a traceability
system that closes all three gaps. The work follows the Design Science Research (DSR)
methodology~\cite{hevner2004}, progressing from gap identification through artefact
construction to empirical verification. The main contributions are:
\begin{itemize}[leftmargin=*, nosep]
  \item A three-layer hybrid storage architecture (blockchain + IPFS + OrbitDB) that stores
        only lightweight hash anchors on-chain, keeping storage costs proportional to
        training scale.
  \item A zero-intrusion Python SDK based on monkey-patching that automatically captures
        training events at Session, Epoch, and Step granularity without modifying user code.
  \item A simplified two-phase atomic commit protocol with compensation-log recovery that
        guarantees idempotent fault recovery across both process-crash and network-partition
        failure modes.
  \item A visual traceability dashboard with a Lineage Map and one-click cryptographic
        verification supporting independent audit by non-technical stakeholders.
  \item Comprehensive empirical evaluation across six quantitative dimensions under
        pre-registered pass/fail thresholds, demonstrating industrial-grade deployability.
\end{itemize}

The remainder of the paper is organised as follows. Section~2 reviews related work and
identifies the three research gaps in detail. Section~3 presents the system architecture and
key design decisions. Section~4 reports empirical evaluation results. Section~5 discusses
broader implications and compares the system against the identified gaps. Section~6
addresses ethical and cross-cultural considerations. Section~7 covers limitations and
future work, followed by conclusions in Section~8.

% ── 2. Background and Related Work ────────────────────────────────────────────
\section{Background and Related Work}

\subsection{Blockchain for AI Auditing}

Soldatos et al.~\cite{soldatos2021} proposed a blockchain-based data provenance framework
for trusted AI, but restricted auditing to the integrity of the input dataset. Zakariae et
al.~\cite{zakariae2025} similarly focused on static dataset verification. Navarro et
al.~\cite{navarro2023} advanced this by re-computing partial training intervals to detect
tampering, yet did not record the dynamic trajectory of parameter changes. Balan et
al.~\cite{balan2025} constructed a rigorous cryptographic framework for end-to-end pipeline
verifiability, but provided no engineering blueprint for blockchain-and-off-chain-storage
collaboration.

\subsection{Hybrid Storage and IPFS Limitations}

Kurada and Baruah~\cite{kurada2023} demonstrated that fully on-chain storage of large-scale
AI metadata is industrially infeasible. Trautwein et al.~\cite{trautwein2022} measured IPFS
at scale and showed that content availability depends on explicit pinning rather than
protocol defaults. Shi et al.~\cite{shi2024} found that a large fraction of IPFS content
requires multiple retrieval attempts and that content providers are highly centralised in a
few cloud nodes. These findings directly informed the mandatory-pin design decision in
our storage architecture.

\subsection{Consistency and Fault Tolerance}

Ruan et al.~\cite{ruan2019} proposed fine-grained on-chain data provenance but restricted
the design to purely on-chain structured data without crash recovery. Belchior et
al.~\cite{belchior2022} identified atomicity and fault tolerance as essential requirements
for cross-layer collaborative systems. Cotroneo et al.~\cite{cotroneo2022} demonstrated
that controlled fault injection is necessary to validate distributed consistency mechanisms.

\subsection{MLOps Tool Limitations}

Idowu et al.~\cite{idowu2024} surveyed 81 ML practitioners and conducted controlled
experiments with 15 developers, finding that while 70\% of respondents used dedicated
tools to track experiments, none of these tools offered cryptographic tamper-proof
guarantees. Historical records can be silently modified by personnel with storage access,
meaning that the integrity of records cannot be independently verified even when an
engineering team is already instrumented with MLOps tooling. This establishes the
distinction between \emph{recording} training metadata and \emph{trusting} it---a gap
that blockchain hash anchoring is uniquely positioned to close.

\subsection{Distributed Tracing and Observability}

Janes et al.~\cite{janes2023} conducted a systematic multivocal review of 30 open-source
distributed tracing tools and found that OpenTelemetry holds significant advantages in
supporting all three observability signal types---traces, metrics, and logs---whereas most
alternatives support only one or two. This peer-reviewed endorsement informed the adoption
of the OpenTelemetry standard in both the Python SDK and the Node.js audit controller,
enabling session-level trace identifiers to serve as cross-service correlation keys across
the full audit pipeline.

% ── 3. System Design ──────────────────────────────────────────────────────────
\section{System Design}

\subsection{Research Methodology}

This project adopts Design Science Research (DSR)~\cite{hevner2004} as its methodological
framework. DSR was selected over pure experimental research (which cannot isolate the
non-determinism of blockchain confirmation times) and case study methods (which describe
existing systems rather than creating new artefacts). The three research questions map
directly to three system artefacts:

\begin{itemize}[leftmargin=*, nosep]
  \item \textbf{RQ1}: How to design a three-layer hybrid storage architecture integrating
        blockchain, IPFS and OrbitDB that provides end-to-end cryptographic verifiability
        while keeping training overhead within an acceptable range?
  \item \textbf{RQ2}: How to design a zero-intrusion metadata collection mechanism that
        automatically captures training data at Session/Epoch/Step granularity and constructs
        an end-to-end verifiable hash chain?
  \item \textbf{RQ3}: How to design an atomic consistency mechanism that maintains
        traceability-chain integrity through automated fault detection and idempotent
        compensation recovery in abnormal scenarios?
\end{itemize}

\subsection{Four-Layer Architecture}

The system is structured into four layers with strict separation of concerns
(Figure~\ref{fig:arch}):

\begin{description}[leftmargin=0pt, nosep]
  \item[\textbf{Layer 1 -- Persistence:}] Blockchain stores only Merkle roots and packet
        hashes; large binary files (checkpoints, model weights) are stored in IPFS with
        mandatory pinning; structured metadata is indexed in OrbitDB for queryable
        off-chain retrieval.
  \item[\textbf{Layer 2 -- Audit Controller:}] A Node.js request handler orchestrates the
        two-phase atomic commit. It anchors hashes on-chain first; off-chain writes proceed
        only after successful anchoring.
  \item[\textbf{Layer 3 -- Data Collection:}] The Python Watcher SDK and Hash Engine
        communicate with Layer 2 solely via HTTP, decoupling the Python training
        environment from the JavaScript storage backend.
  \item[\textbf{Layer 4 -- Presentation:}] A web dashboard (Lineage Map, traffic-light
        status indicators, one-click verification buttons) allows non-technical auditors to
        verify training lineage independently.
\end{description}

\begin{figure}[ht]
  \centering
  \includegraphics[width=0.95\columnwidth]{layers.png}
  \caption{Four-layer system architecture.}
  \label{fig:arch}
\end{figure}

\subsection{Hash Chain Tamper-Proof Mechanism}

Each audit packet contains a \texttt{prev\_hash} field pointing to the hash of its
predecessor, forming an append-only hash chain. Any post-hoc modification to a packet
changes its hash value, causing verification to fail for all subsequent nodes. The
verification endpoint recalculates the current packet hash and compares it against the
\texttt{latestPacketHash} anchored on-chain. IPFS CID tampering is detected by
comparing the stored content-addressing pointer against the expected CID.

\subsection{Non-Intrusive SDK: Monkey-Patch Mechanism}

The proposed solution in the initial design stage described hook-based collection, which
still required users to actively register callbacks. During implementation, this was
replaced by monkey-patching, which dynamically replaces PyTorch internals at runtime
and achieves \emph{zero modification} of user training code. Three patches are applied:

\begin{enumerate}[leftmargin=*, nosep]
  \item \texttt{torch.Tensor.backward} --- automatically extracts scalar loss values on
        each backward pass.
  \item \texttt{torch.optim.Optimizer.\_\_init\_\_} --- injects a step hook for global step
        counting and periodic snapshot triggering.
  \item \texttt{torch.utils.data.DataLoader.\_\_iter\_\_} --- wraps the iterator with a
        sentinel to detect epoch boundaries and trigger checkpoint capture.
\end{enumerate}

Initialisation requires a single import and one line of user code; all subsequent training
code remains unchanged. The trade-off is tight coupling to PyTorch; multi-framework
support is reserved as future work.

\subsection{Software Engineering Principles}

Four established principles are systematically applied. \textbf{Separation of Concerns}
is enforced across all layers: migrating the blockchain backend from Ethereum to
Hyperledger Fabric requires changes only in \texttt{blockchain.js}. \textbf{Programming
to Interface} ensures all SDK interactions with storage backends occur via HTTP endpoints
configured through environment variables, leaving the SDK agnostic to the underlying
store. \textbf{Agile Development} (five two-week sprints) enabled the monkey-patch
correction to be detected and applied during Sprint~2 without disrupting the research
goals. \textbf{Observability by Design} embeds the three OpenTelemetry pillars (metrics,
logs, traces) from the outset~\cite{janes2023}: Prometheus exporters continuously stream
performance counters, structured JSON logs record every critical operation, and distributed
trace IDs link SDK collection events to their corresponding on-chain anchoring
transactions.

\subsection{Simplified Two-Phase Atomic Commit}

The atomic consistency mechanism follows a strict ordering:

\begin{enumerate}[leftmargin=*, nosep]
  \item \textbf{Phase 1 (On-chain anchoring):} Submit the packet hash to the blockchain.
        If this fails, return HTTP 500 immediately---no off-chain write is performed.
  \item \textbf{Phase 2 (Off-chain write):} Write to OrbitDB only after Phase 1 succeeds.
        If Phase 2 fails, the transaction hash is appended to \texttt{orphan-txs.log}.
\end{enumerate}

A file-system watcher (\texttt{fs.watch}) monitors \texttt{orphan-txs.log} and triggers
\texttt{recoverOrphanTxs()} automatically after a 2-second debounce. Idempotency is
guaranteed by checking \texttt{db.get(entryId)} before each write: if the record already
exists from a prior recovery attempt, the write is skipped.

% ── 4. Evaluation ─────────────────────────────────────────────────────────────
\section{Empirical Evaluation}

Evaluation followed the DSR principle that \emph{criteria precede implementation}. Five
dimensions and eight quantifiable indicators with explicit pass thresholds were defined
before Sprint 5 began. A dual-path verification logic was adopted: (i) controlled
experiments with paired-sample design for RQ1; (ii) cryptographic verification and
adversarial fault injection for RQ2 and RQ3.

\subsection{Dimension 1: Training Performance Overhead (RQ1)}

Five repeated runs of an MNIST training task (10 epochs, batch size 64, step interval 200)
were performed under audit-disabled and audit-enabled conditions. Results:

\begin{table}[ht]
  \centering
  \caption{Training time: control vs.\ audit group ($n=5$).}
  \label{tab:perf}
  \small
  \begin{tabular}{lcc}
    \toprule
    Condition & Mean time (s) & Std (s) \\
    \midrule
    Control (no audit) & 193.5 & 1.8 \\
    Audit enabled       & 199.9 & 3.8 \\
    \midrule
    Overhead            & \multicolumn{2}{c}{$+3.28\%$ \quad (\textbf{PASS} $< 5\%$)} \\
    \textit{t}-test     & \multicolumn{2}{c}{$t = -3.337$, $p = 0.0103$} \\
    \bottomrule
  \end{tabular}
\end{table}

The higher variance in the audit group ($\pm3.8$\,s vs.\ $\pm1.8$\,s) reflects the
concentration of overhead in session-level one-time operations (startSession transaction,
IPFS upload of dataset and model), rather than uniform per-epoch accumulation. At the
epoch level the overhead was $+2.54\%$ and throughput reduction was $2.48\%$
(49.17 $\to$ 47.95 batches/s). Prometheus queue-depth monitoring confirmed zero request
backlog throughout the experiment.

\subsection{Dimension 2: Audit Effectiveness (RQ2)}

Six targeted tampering scenarios were injected programmatically via the
\texttt{/debug/tamper} interface, covering step loss metrics, epoch checkpoint hash and
CID, and session dataset/model hash and CID. All six scenarios triggered an
\texttt{invalid} alert; the detection rate was \textbf{100\%} with zero missed detections
(\textbf{PASS}).

\subsection{Dimension 3: Verification Response Time (RQ2)}

Verification latency was measured via a Prometheus histogram metric
(\texttt{audit\_verify\_duration\_ms}) across four verification types (session, dataset,
model, checkpoint). All four types achieved a P99 latency below 500\,ms, well within
the 2000\,ms SLA target (\textbf{PASS}).

\subsection{Dimension 4: Storage Scalability (RQ1)}

Off-chain write events were fitted with linear regression. The result was:
\[
  \text{entries} = 17.84t + 0.99, \quad R^2 = 0.99985 \quad (\textbf{PASS} > 0.99).
\]
On-chain gas consumption showed a corresponding linear growth pattern, with an
average cost of 31,988 gas per packet. The \texttt{syncProgress} method dominated
gas consumption (94.7\%), consistent with its role as the high-frequency per-packet
operation. These results confirm cost predictability for industrial-scale deployment.

\subsection{Dimension 5: Atomic Consistency and Fault Recovery (RQ3)}

Following the fault injection methodology of Cotroneo et al.~\cite{cotroneo2022}, two
fault types were injected between successful on-chain anchoring and the OrbitDB write:
\texttt{process\_crash} (immediate exception) and \texttt{network\_partition} (10-second
timeout). Each type was repeated three times (6 rounds total). Results:

\begin{table}[ht]
  \centering
  \caption{Fault recovery and idempotency results ($n=6$ rounds).}
  \label{tab:fault}
  \small
  \begin{tabular}{lcc}
    \toprule
    Metric & Result & Target \\
    \midrule
    Orphan transactions recovered  & 6/6 & 6/6 \\
    Orphans remaining after recovery & 0 & 0 \\
    Duplicate records (idempotency) & 0 & 0 \\
    \midrule
    Recovery rate & \textbf{100\%} & 100\% \\
    Idempotency violations & \textbf{0} & 0 \\
    \bottomrule
  \end{tabular}
\end{table}

Both fault types are handled uniformly through the \texttt{orphan-txs.log} compensation
path, decoupling the recovery logic from the specific failure mode---a design decision
aligned with the fault-tolerance requirements of Belchior et al.~\cite{belchior2022}.

\subsection{Dimension 6: Audit Capture Coverage (RQ2)}

The \texttt{/audit/session/:id/coverage} endpoint reported 47/47 step packets and 10/10
epoch packets captured (both coverage ratios = 1.0000, \textbf{PASS}). Combined with
the 100\% tamper detection result, this establishes that the SDK not only collects all
data completely but also ensures the collected data is cryptographically unforgeable.

\subsection{Summary}

\begin{table}[ht]
  \centering
  \caption{Evaluation summary: all eight indicators pass.}
  \label{tab:summary}
  \small
  \begin{tabular}{lccc}
    \toprule
    Criterion & Target & Measured & Result \\
    \midrule
    Training overhead        & $< 5\%$       & $+3.28\%$       & PASS \\
    Audit effectiveness      & 100\%         & 100\%           & PASS \\
    Verification response    & $< 2000$\,ms  & P99 498\,ms     & PASS \\
    Storage linear growth    & $R^2 > 0.99$  & $0.99985$       & PASS \\
    Fault recovery rate      & 100\%         & 100\%           & PASS \\
    Idempotency (duplicates) & 0             & 0               & PASS \\
    Step capture coverage    & 1.0 (100\%)   & 100.00\%        & PASS \\
    Epoch capture coverage   & 1.0 (100\%)   & 100.00\%        & PASS \\
    \bottomrule
  \end{tabular}
\end{table}

% ── 5. Discussion ─────────────────────────────────────────────────────────────
\section{Discussion}

\subsection{Closing the Three Research Gaps}

The evaluation results confirm that all three identified gaps have been closed at the
engineering level. Gap~1 (theory-to-engineering) is addressed by the three-layer hybrid
architecture, which was deployed, instrumented, and empirically measured rather than
remaining a theoretical proposal. Linearly growing storage cost directly refutes the
scalability concern of Kurada and Baruah~\cite{kurada2023}. Gap~2 (dynamic audit
granularity) is addressed by the monkey-patch SDK achieving complete capture coverage
at both Step and Epoch levels, with every record rendered cryptographically unforgeable
through hash-chain linkage. This extends the tamper-proof guarantee from the final model
artefact---the limit of prior work~\cite{navarro2023}---to the complete dynamic training
trajectory. Gap~3 (atomic consistency) is resolved by the two-phase commit protocol,
whose fault recovery was verified empirically via controlled fault injection rather than
static design argument alone, fulfilling the methodological requirement of Cotroneo et
al.~\cite{cotroneo2022}.

\subsection{Comparison with Prior Work and Scientific Reproducibility}

Compared to Navarro et al.~\cite{navarro2023}, the present system provides broader
coverage across three orthogonal dimensions: collection granularity (Session/Epoch/Step
vs.\ final-model only), storage architecture (hybrid with atomic fault recovery vs.\
on-chain only), and verification scope (dataset, checkpoint, model, and lineage
independently vs.\ model integrity only). Beyond regulatory compliance, the system
has direct applicability to the AI reproducibility crisis: by anchoring step-level
hyperparameters and loss trajectories to an immutable chain, third-party researchers
can cryptographically verify claimed training configurations without access to the
original infrastructure---shifting reproducibility from a social trust problem to a
technical one~\cite{idowu2024}.

% ── 6. Ethical and Cultural Considerations ────────────────────────────────────
\section{Ethical and Cultural Considerations}

\subsection{Privacy vs.\ Cryptographic Auditability}

While the Merkle tree structure does not expose raw data, it introduces a membership
inference vulnerability~\cite{shokri2017}: an attacker possessing a data sample can
request a Merkle proof path to confirm whether that sample belongs to the training set
without accessing the original dataset. Countermeasures include role-based access
control on proof-path queries, salted leaf-node hashes to prevent external replication, and
optional homomorphic encryption for sensitive metadata fields.

\subsection{Right to Be Forgotten vs.\ Immutable Blockchain}

GDPR Article 17 grants data subjects the right of erasure. Once a personal-data hash is
anchored on an immutable chain, the hash record persists even after the original data is
deleted~\cite{politou2018}. A hash revocation mechanism---publishing an on-chain
revocation statement to mark a specific hash as invalid---provides a technical
approximation of erasure without modifying immutable records. For EU deployments,
consortium or private chain architectures with agreed modification rules are preferable.

\subsection{Accountability vs.\ Responsibility Evasion}

The traceability system may be misused to attribute bias outcomes to ``uncontrollable
data distribution''~\cite{pasquale2015}, shielding organisations from moral
accountability. System documentation must explicitly state that complete training records
are a \emph{necessary but not sufficient} condition for accountability: the organisation
retains full moral and legal responsibility for training data representativeness, label
quality, and deployment boundary decisions.

\subsection{Cross-Cultural Regulatory Compliance}

The interface-oriented programming principle enables compliance configuration to be
switched at the interface level without modifying core SDK logic: private IPFS clusters
with geo-fencing for EU GDPR; fully private deployment for Chinese data localisation
requirements (Data Security Law, PIPL); modular state-level configuration for US markets.
The separation-of-concerns architecture additionally supports flexible switching among
public chains (high transparency), consortium chains (industry self-regulation), and
private chains (internal audit), adapting to different cultural trust models without
restructuring the system core.

% ── 7. Limitations and Future Work ────────────────────────────────────────────
\section{Limitations and Future Work}

The prototype carries four limitations. \textbf{Framework coupling:} the Watcher SDK is
tightly bound to PyTorch; adapting to TensorFlow or JAX requires re-implementing the
three monkey-patch targets for each framework's internals. \textbf{IPFS availability:}
single-node pinning poses availability risks in production~\cite{trautwein2022,shi2024};
cross-node pin replication or Filecoin integration is needed. \textbf{Workload scope:}
experiments used MNIST as a proxy; transformer architectures and multi-GPU distributed
training may exhibit different overhead profiles, particularly for the model-file IPFS
upload step. \textbf{Fault coverage:} only process-crash and network-partition modes were
tested; Byzantine faults and concurrent write conflicts remain unverified.

Future directions include: multi-framework SDK support via a framework-agnostic
abstraction layer; multi-node IPFS pin management; SaaS deployment to lower adoption
barriers for SMEs; differential privacy noise injection into evidential metadata to
mitigate membership inference risks; and red-team testing under a responsible disclosure
policy following open-source release.

% ── 8. Conclusion ─────────────────────────────────────────────────────────────
\section{Conclusion}

This paper presented a blockchain-enabled AI training traceability system designed and
evaluated under the DSR methodology. Three research gaps---the theory-to-engineering
chasm in hybrid storage architecture, the insufficient granularity of dynamic process
auditing, and the systematic absence of atomic consistency guarantees---were each
addressed by a corresponding system artefact: the four-layer blockchain/IPFS/OrbitDB
architecture, the monkey-patch Watcher SDK, and the two-phase atomic commit protocol
with compensation-log recovery. All eight pre-registered evaluation indicators were
passed in controlled experiments, confirming industrial-grade performance overhead,
complete tamper detectability, fast one-click verification, predictable storage cost
growth, and full fault recoverability with zero idempotency violations.

The system advances the state of the art beyond prior engineering implementations
by extending tamper-proof guarantees to the full dynamic training lifecycle rather than
the final model alone, and by providing atomic consistency under real fault conditions
rather than by design assumption. Ethical analysis identified residual tensions around
privacy, data sovereignty, and accountability that cannot be resolved through technical
means alone, but that can be substantially mitigated through architectural choices,
access control policies, and explicit governance documentation. The result is a
deployable, auditable, and compliance-ready foundation for trustworthy AI governance.

% ── References ────────────────────────────────────────────────────────────────
\balance
{\small
\setlength{\itemsep}{0pt}
\setlength{\parsep}{0pt}
\begin{thebibliography}{99}

\bibitem{balan2025}
K.~Balan, R.~Learney, and T.~Wood,
``A Framework for Cryptographic Verifiability of End-to-End AI Pipelines,''
\textit{Proc.\ 10th ACM Int'l Workshop on Security and Privacy Analytics (IWSPA '25)},
pp.\ 49--59, 2025.

\bibitem{belchior2022}
R.~Belchior, A.~Vasconcelos, M.~Correia, and T.~Hardjono,
``Hermes: Fault-tolerant middleware for blockchain interoperability,''
\textit{Future Generation Computer Systems}, vol.~129, pp.\ 236--251, 2022.

\bibitem{cotroneo2022}
D.~Cotroneo, L.~De~Simone, P.~Liguori, and R.~Natella,
``Fault Injection Analytics: A Novel Approach to Discover Failure Modes in Cloud-Computing Systems,''
\textit{IEEE Trans.\ Dependable and Secure Computing}, vol.~19, no.~3, pp.\ 1476--1491, 2022.

\bibitem{hevner2004}
A.~R.~Hevner, S.~T.~March, J.~Park, and S.~Ram,
``Design Science in Information Systems Research,''
\textit{MIS Quarterly}, vol.~28, no.~1, pp.\ 75--105, 2004.

\bibitem{idowu2024}
S.~Idowu, O.~Osman, D.~Str\"{u}ber, and T.~Berger,
``Machine learning experiment management tools: A mixed-methods empirical study,''
\textit{Empirical Software Engineering}, vol.~29, no.~4, p.~74, 2024.

\bibitem{janes2023}
A.~Janes, X.~Li, and V.~Lenarduzzi,
``Open tracing tools: Overview and critical comparison,''
\textit{Journal of Systems and Software}, vol.~204, p.~111793, 2023.

\bibitem{kurada2023}
V.~R.~Kurada and P.~K.~Baruah,
``Blockchain Enabled, Collaborative Platform for AI as a Service,''
\textit{ICTACT Journal on Soft Computing}, vol.~13, no.~3, pp.\ 2909--2916, 2023.

\bibitem{navarro2023}
E.~Navarro, K.~J.~Standing, G.~G.~Dagher, and T.~Andersen,
``Ensuring Trustworthy Neural Network Training via Blockchain,''
\textit{Proc.\ IEEE 5th Int'l Conf.\ Cognitive Machine Intelligence (CogMI)},
pp.\ 31--40, 2023.

\bibitem{pasquale2015}
F.~Pasquale,
\textit{The Black Box Society: The Secret Algorithms That Control Money and Information}.
Harvard University Press, 2015.

\bibitem{politou2018}
E.~Politou, E.~Alepis, and C.~Patsakis,
``Forgetting personal data and revoking consent under the GDPR: Challenges and proposed solutions,''
\textit{Journal of Cybersecurity}, vol.~4, no.~1, p.~tyy001, 2018.

\bibitem{ruan2019}
P.~Ruan et al.,
``Fine-grained, secure and efficient data provenance on blockchain systems,''
\textit{Proc.\ VLDB Endowment}, vol.~12, no.~9, pp.\ 975--988, 2019.

\bibitem{sciammarelli2026}
J.~Sciammarelli,
``Blockchain-Enabled Infrastructure for Secure and Auditable AI Systems,''
\textit{REMUNOM}, vol.~13, no.~04, pp.\ 1--13, 2026.

\bibitem{shi2024}
R.~Shi et al.,
``A Closer Look into IPFS: Accessibility, Content, and Performance,''
\textit{Proc.\ ACM Meas.\ Anal.\ Comput.\ Syst.}, vol.~8, no.~2, pp.\ 20:1--20:31, 2024.

\bibitem{shokri2017}
R.~Shokri, M.~Stronati, C.~Song, and V.~Shmatikov,
``Membership Inference Attacks Against Machine Learning Models,''
\textit{Proc.\ IEEE Symp.\ Security and Privacy (SP)}, pp.\ 3--18, 2017.

\bibitem{soldatos2021}
J.~Soldatos, A.-M.~Despotopoulou, N.~Kefalakis, and B.~Ipektsidis,
``Blockchain based Data Provenance for Trusted Artificial Intelligence,''
in \textit{Becoming a Platform in Europe}, Now Publishers, 2021.

\bibitem{trautwein2022}
D.~Trautwein et al.,
``Design and Evaluation of IPFS: A Storage Layer for the Decentralised Web,''
\textit{Proc.\ ACM SIGCOMM}, pp.\ 739--752, 2022.

\bibitem{zakariae2025}
S.~Zakariae, A.~Ouidad, and E.~B.~E.~I.~Younes,
``Blockchain-Assisted Assurance of Data Integrity in AI Model Training,''
\textit{Informatica}, vol.~49, no.~18, 2025.

\end{thebibliography}
}

\end{document}
